% ===================================================================
%  TABELO N° 1 — La skolo, la liceo, la docochambro.
%  Feuillets 7 a 12 du fac-simile = folios 5 a 10.
%  Correspondance : folio imprime = numero de feuillet - 2.
%
%  Les marques « %%K » ne sont pas des commentaires ordinaires : elles
%  portent la CLE D'ALIGNEMENT qui apparie, dans la page de lecture, ce
%  bloc-ci et son homologue francais. La cle se lit
%      t<tableau>-<numero d'alinea numerote>-<rang dans l'alinea>
%  et ne depend d'aucune pagination : les deux editions ne coupent pas
%  leurs pages aux memes endroits, mais elles numerotent les memes
%  alineas — c'est le seul point d'ancrage que les deux livrets
%  partagent, et il est de l'auteur, non du transcripteur.
% ===================================================================

% --- Feuillet 7 (folio 5, non imprime : page d'ouverture) ----------
\begin{VUpage}[7]{}
%%K t01-apar-1 apar
\VUsaut{2.0mm}
\VUcentre{12.6pt}{120}{EXPLIKO - LIBRETO}
\VUsaut{3.2mm}
\VUcentre{8.4pt}{160}{DI}
\VUsaut{2.6mm}
\VUtitre{15.6pt}{0}{la Delmas - tabeli helpanta}
\VUsaut{3.4mm}
\VUornamento{9pt}
\VUsaut{5.0mm}
\VUcentre{13.2pt}{90}{UNESMA SERIO}
\VUsaut{3.0mm}
\VUfilet{16mm}
\VUsaut{5.6mm}
\VUtitre{15.0pt}{60}{TABELO N\textsuperscript{o} 1}
\VUsaut{1.4mm}
\VUtitre{11.4pt}{0}{La skolo, la liceo, la docochambro.}
\VUsaut{2.2mm}
\VUfilet{20mm}
\VUsaut{6.4mm}

%%K t01-tit-1 sub
\VUpk{10.2pt}{40}{Deskripto generala.}
\VUsaut{3.0mm}

%%K t01-01-1 p
\VUblancAlinea
1. --- Ica tabelo reprezentas \VUgras{docochambro}\nl
en \VUgras{liceo.} En ica docochambro esas ok \VUgras{tabli} \textsuperscript{(18)}\nl
ed ok \VUgras{benki} \textsuperscript{(32)}. Sur singla benko sidas tri ler\cc
nanti. La \VUgras{docisto} \textsuperscript{(7)} sidas sur \VUgras{stulo} \textsuperscript{(8)} en sua\nl
\VUgras{katedro} \textsuperscript{(6)}.

%%K t01-02-1 p
\VUblancAlinea
2. --- Sinistre di la katedro trovesas \VUgras{ar}\cc
\VUgras{moro} \textsuperscript{(15)} \VUgras{vitrizita.} Dextre esas la \VUgras{nigra ta}\cc
\VUgras{belo} \textsuperscript{(1)} kun la \VUgras{vishilo} \textsuperscript{(2)} e la \VUgras{kreto} \textsuperscript{(3)}.

%%K t01-02-2 p
\VUblancAlinea
Super la katedro pendas \VUgras{mural tabeli} \textsuperscript{(9, 11,}\nl
\textsuperscript{12)} e \VUgras{mapo} \textsuperscript{(10)}.

%%K t01-03-1 p
\VUblancAlinea
3. ---Ne omna lernanti sidas en sua \VUgras{plaso} \textsuperscript{(17)}.\nl
Ioannes \textsuperscript{(4)} stacas an la tabelo, il tenas \VUgras{la}\nl
vishilo ed efacas la \VUgras{geometrial figuri.}

%%K t01-03-2 p
\VUblancAlinea
Petrus esas proxim la katedro, il kolektas la\nl
\VUgras{skribo-taski} \textsuperscript{(5)}, ed adportas li al docisto.
\end{VUpage}

% --- Feuillet 8 (folio 6) ------------------------------------------
\begin{VUpage}[8]{6}
%%K t01-04-1 p
\VUblancAlinea
4. --- \VUgras{Ica} docisto esas la docisto pri la \VUgras{vi}\cc
\VUgras{vanta lingui.} Il docas a la lernanti parolar, lek\cc
tar e skribar stranjera lingui \VUgras{(Germana, An}\cc
\VUgras{gla, Hispana, Italiana, Rusa} o \VUgras{Franca)}.

%%K t01-05-1 p
\VUblancAlinea
5. --- La docochambro esas tre vasta et tre\nl
klara. En la fundo esas larja \VUgras{fenestro} kun du\nl
\VUgras{postigi} \textsuperscript{(78)} e la \VUgras{pordo vitrizita} por aerizar e lu\cc
mizar olu.

%%K t01-05-2 p
\VUblancAlinea
Ica docochambro ne esas kolda. En la mezo\nl
esas por varmigar ol tre bona \VUgras{furnelo} \textsuperscript{(46)} kun\nl
lua \VUgras{tubo} \textsuperscript{(45)}.

%%K t01-05-3 p
\VUblancAlinea
An la parieto esas fixigita \VUgras{vest-hoki} \textsuperscript{(58)}, de\nl
li pendas la kapvesti e la manteli di la skolani.

%%K t01-tit-2 sub
\VUsaut{5.2mm}
\VUpk{10.2pt}{40}{La ludo-korto}
\VUsaut{3.6mm}

%%K t01-06-1 p
\VUblancAlinea
6. --- La \VUgras{horlojo} \textsuperscript{(63)} indikas non kloki e kin\nl
minuti.

%%K t01-06-2 p
\VUblancAlinea
La \VUgras{amuz-tempo} \textsuperscript{(76)} jus finis e la leciono riko\cc
mencas. L'amuzeyo situesas en la \VUgras{korto} \textsuperscript{(75)}. Ni\nl
vidas kelka lernanti qui ludas. \VUgras{Submaestro} \textsuperscript{(74)}\nl
surveyas li.

%%K t01-07-1 p
\VUblancAlinea
7. --- En la korto ni pluse vidas diversa per\cc
soni, nome l'\VUgras{inspektisto} \textsuperscript{(73)}, la skolestro \VUgras{(li}\cc
\VUgras{ceestro)} \textsuperscript{(71)} e la \VUgras{censoro} \textsuperscript{(72)}.

%%K t01-07-2 p
\VUblancAlinea
L'inspektisto inspektas la skoli, la liceestro\nl
direktas la liceo, la censoro surveyas la studii.

%%K t01-07-3 p
\VUblancAlinea
La quaresma persono esas la \VUgras{pordisto} \textsuperscript{(77)}. Il\nl
portas sub brakio registro por enskribar l'ab\cc
senti.

%%K t01-08-1 p
\VUblancAlinea
8. --- En la fundo di la korto esas la \VUgras{gimnas}\cc
\VUgras{tikal aparati} \textsuperscript{(70)} e la laboreyo por la \VUgras{manu}\cc
\VUgras{labori} \textsuperscript{(69)}.

%%K t01-08-2 p
\VUblancAlinea
Dextre dil tabelo ni videskas la \VUgras{doco-cham}\cc
\VUgras{bri} pri \VUgras{muziko} e \VUgras{desegno}.
\end{VUpage}

% --- Feuillet 9 (folio 7) ------------------------------------------
%  Cette page porte la SEULE note du tableau : l'appel est un
%  asterisque entre parentheses, et la note explique la latinisation
%  des prenoms. Elle est propre a l'edition ido — le livret francais ne
%  la connait pas, et la page de lecture la donne donc sans vis-a-vis.
\begin{VUpage}[9]{7}
%%K t01-noto-1 noto
\VUnotes{143.2mm}{%
(*) Por la \textit{baptonomi} on konsilas anke transskribar li\nl
segun la formo Latina. To esas l'unika moyeno eskapar di\cc
vergi sennombra. Cetere quale selektar inter \textit{Giovanni, Jean,}\nl
\textit{Johann, Jan, Hans, John, Ivan,} e. c.? Ka ni kreos vortaro\nl
specala por la baptonomi? Ni cherpez de la Latina lia no\cc
minativo e dicez : \VUgras{Ioannes, Ioanna; Iulius, Iulia; Iakobus;}\nl
\VUgras{Andreas; Lukas.} (Gram. Kompleta).%
}

%%K t01-tit-3 sub
\VUpk{10.2pt}{40}{Detaloza deskripto.}
\VUsaut{2.4mm}

%%K t01-tit-4 sub
\VUpk{10.2pt}{40}{La bona lernanti.}
\VUsaut{3.0mm}

%%K t01-09-1 p
\VUblancAlinea
9. --- La lernanti dil unesma benko dextre\nl
dil katedro esas \VUgras{Henrikus} (*), \VUgras{Stefanus} e \VUgras{Ka}\cc
\VUgras{rolus.}

%%K t01-09-2 p
\VUblancAlinea
Henrikus esas tre atencoza e laborema, il\nl
askoltas la docisto.

\VUblancAlinea
Sur sua tablo il havas \VUgras{kayero} \textsuperscript{(28)}, \VUgras{libro} \textsuperscript{(29)},\nl
\VUgras{lexiko} \textsuperscript{(30)} e \VUgras{skribiluyo} \textsuperscript{(31)}. Lua libro havas\nl
multa \VUgras{pagini}, singla chapitro kontenas plura\nl
\VUgras{paragrafi} e singla \VUgras{lineo} multa \VUgras{vorti.}

%%K t01-09-4 p
\VUblancAlinea
Lua vicino Stefanus \textsuperscript{(24)} esas fervoroza. Il\nl
notas sur sua kayero la leciono por la sequonta\nl
foyo. Il havas bela \VUgras{skribo-formo.} Lua libro\nl
esas klozita. Ni vidas nur la \VUgras{kovrilo} \textsuperscript{(27)}. Proxim\nl
lu trovesas lua \VUgras{kompaso} \textsuperscript{(26)}.

%%K t01-09-5 p
\VUblancAlinea
Karolus \textsuperscript{(23)} tenas per la manuo sua libro,\nl
il apertas olu por rilektar sua leciono. Il havas,\nl
quale singla lernanto, \VUgras{inkuyo} \textsuperscript{(25)} avan su. Il\nl
esas obediera.

%%K t01-10-1 p
\VUblancAlinea
10. --- An la tablo apuda esas \VUgras{Ludovikus,}\nl
\VUgras{Antonius} e \VUgras{Paulus.} Ludovikus recitas sua \VUgras{le}\cc
\VUgras{ciono} \textsuperscript{(22)} e respondizas la \VUgras{questioni} dil docisto.\nl
Antonius \textsuperscript{(21)} esas tre neatencoza, mem kelka\cc
foye la docisto punisas lu. Paulus \textsuperscript{(20)} esas bona\nl
kamarado, afabla et komplezema. Il esas tre\nl
atencoza.
\end{VUpage}

% --- Feuillet 10 (folio 8) -----------------------------------------
\begin{VUpage}[10]{8}
%%K t01-tit-5 sub
\VUsaut{4.2mm}
\VUpk{10.2pt}{40}{La mala lernanti.}
\VUsaut{3.2mm}

%%K t01-11-1 p
\VUblancAlinea
11. --- Dop Henrikus trovesas \VUgras{Georgius} \textsuperscript{(38)}.\nl
Il ne esas mala kerlo, ma il esas \VUgras{babilema,}\nl
neatencoza e petulema. Lua \VUgras{frivoleso} e \VUgras{desobe}\cc
\VUgras{diemeso} efektigas \VUgras{trublo} dum la leciono ed ofte\nl
il meritas sever \VUgras{averto.} Lua \VUgras{klado-kayero} \textsuperscript{(35)}\nl
esas sordida, il \VUgras{inko-makulizas} lu per sua\nl
\VUgras{plumo} \textsuperscript{(34)}, pro to, il sempre uzas sua \VUgras{efaco}\cc
\VUgras{kauchuko} \textsuperscript{(36)}; lua \VUgras{kayeruyo} \textsuperscript{(33)} esas lacerita,\nl
nam il esas tre neglijera. Pro quo il kunportas\nl
sua \VUgras{krayoniero} \textsuperscript{(37)} en la docochambro di \VUgras{la}\nl
vivanta lingui ?

%%K t01-11-2 p
\VUblancAlinea
Lua vicino Edmundus \textsuperscript{(39)} ne amas lu. Il esas\nl
advere kelke indolenta, ma il esas tacera e\nl
tranquila. Il ne babilas; ma, vice askoltar, il\nl
preferas lektar la \VUgras{rakonti} di sua \VUgras{lekto-libro.}

%%K t01-11-3 p
\VUblancAlinea
La triesma lernanto di ca benko esas \VUgras{Ludo}\cc
\VUgras{vikus} \textsuperscript{(40)}. Il esas diligenta. Il skribas sur\nl
blanka \VUgras{paperfolio.} Avan su, il pozis sua \VUgras{absor}\cc
\VUgras{biva papero} \textsuperscript{(41)}.

%%K t01-12-1 p
\VUblancAlinea
12. --- La lernanti dil duesma benko, sinistre\nl
dil docisto, esas tre yuna. L'unesma, la mikra\nl
\VUgras{Emilius,} ne ja savas tre bone skribar, il kun\cc
portas sua \VUgras{ardezo} \textsuperscript{(42)}, sua \VUgras{ardeza krayono} e sua\nl
\VUgras{sponjo} \textsuperscript{(43)}.

%%K t01-12-2 p
\VUblancAlinea
Apud ilu esas \VUgras{Augustus} e \VUgras{Bernardus} \textsuperscript{(44)}. Ica\nl
bone laboras, ma lua \VUgras{vestizeso} esas tre negli\cc
jita.

%%K t01-13-1 p
\VUblancAlinea
13. --- Dop li trovesas tri \VUgras{indolenti. Alber}\cc
\VUgras{tus} ne lernas sua lecioni. \VUgras{Iakobus} \textsuperscript{(47)} dormas\nl
vice askoltar. Lua \VUgras{indolenteso} efektigas a \VUgras{lu}\nl
\VUgras{reprochi} e mem \VUgras{punisi.} Fine \VUgras{Kristianus} \textsuperscript{(48)} esas
\parplein
\end{VUpage}

% --- Feuillet 11 (folio 9) -----------------------------------------
\begin{VUpage}[11]{9}
%%K t01-13-1 p suite
\VUcontinue
tro neserioza. Il male \VUgras{kondutas} e nule esforcas\nl
por emendar su.

%%K t01-14-1 p
\VUblancAlinea
14. --- Lua vicini, kontraste, esas bonkon\cc
duta. \VUgras{Ernestus} \textsuperscript{(51)} prizas l'ordino e la regulo\cc
zeso, il sempre uzas sua \VUgras{linealo} \textsuperscript{(50)} e sua\nl
\VUgras{krayono} \textsuperscript{(49)}. Il esas \VUgras{externo.}

%%K t01-14-2 p
\VUblancAlinea
\VUgras{Franciskus} \textsuperscript{(52)} esas \VUgras{interno}, il esas vestizita\nl
per uniformo, manjas e kushas su en la liceo.\nl
Il esas bona \VUgras{kamarado.}

%%K t01-14-3 p
\VUblancAlinea
Mauricius talias sua \VUgras{krayono} per sua \VUgras{kulte}\cc
\VUgras{leto} \textsuperscript{(56)}, il esas tre sorgema.

%%K t01-14-4 p
\VUblancAlinea
Il kunportas omna sua libri, sua \VUgras{grama}\cc
\VUgras{tiko} \textsuperscript{(55)} sua \VUgras{notlibreto} \textsuperscript{(54)}, e mem \VUgras{skribkuse}\cc
\VUgras{no} \textsuperscript{(53)}. Lua \VUgras{skolsako} \textsuperscript{(57)} esas an la sulo dop\nl
ilu.

%%K t01-14-5 p
\VUblancAlinea
La du tabli funde di la docochambro oku\cc
pesas da \VUgras{bona lernanti} \textsuperscript{(19)}.

%%K t01-tit-6 sub
\VUsaut{4.6mm}
\VUpk{10.2pt}{40}{Desegno e muziko.}
\VUsaut{3.4mm}

%%K t01-15-1 p
\VUblancAlinea
15. --- En la \VUgras{desegnochambro} esas gracioza\nl
\VUgras{modeli} akrochita an la muro e \VUgras{lampo} \textsuperscript{(62)} kun\nl
\VUgras{paralumo}, suspendita ye la plafono. La \VUgras{do}\cc
\VUgras{cisto} \textsuperscript{(60)} esas tre pacienta, il korektigas la \VUgras{de}\cc
\VUgras{segnuri} \textsuperscript{(61)} di la lernanti e docas a li desegnar\nl
\VUgras{busto} \textsuperscript{(59)} segun la naturo.

%%K t01-16-1 p
\VUblancAlinea
16. --- La \VUgras{muzikdoceyo} semblas tre intere\cc
siva. En ol on lernas lektar la \VUgras{muziko}, solfejar\nl
la \VUgras{gamnoti} \textsuperscript{(67)} skribita sur la \VUgras{linearo} (do, re,\nl
mi, fa, sol, la, si\,=\,C. D. E. F. G. A. B.) kono\cc
car la \VUgras{klefi} e kantar charmanta \VUgras{melodii.} On\nl
pozas la muzikaji sur la \VUgras{pupitro} \textsuperscript{(65)}. Un de la\nl
lernanti \VUgras{pianopleas} \textsuperscript{(64)} e la \VUgras{muzik-docisto} \textsuperscript{(66)}\nl
\VUgras{skandas} \textsuperscript{(68)}.
\end{VUpage}

% --- Feuillet 12 (folio 10) ----------------------------------------
\begin{VUpage}[12]{10}
%%K t01-17-1 p
\VUblancAlinea
17. --- Ni videskas angulo di la \VUgras{laboreyo} por\nl
\VUgras{manulabori} \textsuperscript{(69)}, ube la pueri darfas lernar ra\cc
botar la ligno e limar la fero. To ne existas en\nl
omna licei.

%%K t01-tit-7 sub
\VUsaut{5.0mm}
\VUpk{10.2pt}{40}{La doceyo pri la cienci.}
\VUsaut{3.6mm}

%%K t01-18-1 p
\VUblancAlinea
18. --- La docisto pri matematiko docas \VUgras{a}\nl
la lernanti \VUgras{geometrio, aritmetiko, kalkulo} e la\nl
\VUgras{metral sistemo}. Ni vidas sur la tabelo la pre\cc
cipua figuri geometriala : \VUgras{rondo} \textsuperscript{(a)}, \VUgras{quadra}\cc
\VUgras{to} \textsuperscript{(b)}, \VUgras{triangulo} \textsuperscript{(c)}, \VUgras{lineo} \textsuperscript{(d)}.

%%K t01-18-2 p
\VUblancAlinea
Ni anke vidas un \VUgras{operaco}, ol ne esas \VUgras{adi}\cc
\VUgras{ciono}, nek \VUgras{sustraciono}, nek \VUgras{multipliko}, ol ya\nl
esas \VUgras{divido} \textsuperscript{(e)}.

%%K t01-19-1 p
\VUblancAlinea
19. --- Sur la tabelo dil metral sistemo, ni\nl
vidas : \VUgras{metro} \textsuperscript{(a)}, \VUgras{balanco} \textsuperscript{(b)} e lua \VUgras{ponderili,}\nl
\VUgras{litro} \textsuperscript{(e)}, \VUgras{dekalitro} \textsuperscript{(f)}, \VUgras{decilitro} e \VUgras{metral kubo} \textsuperscript{(d)}.

%%K t01-19-2 p
\VUblancAlinea
La lernanti anke studias \VUgras{naturcienco} \textsuperscript{(11)},\nl
zoologio \textsuperscript{(a)}, botaniko \textsuperscript{(b)}, geologio \textsuperscript{(c)} e \VUgras{histo}\cc
\VUgras{rio} \textsuperscript{(12)}.

%%K t01-19-3 p
\VUblancAlinea
En l'armoro esas la aparati pri \VUgras{fiziko} \textsuperscript{(15)} e\nl
pri \VUgras{kemio} \textsuperscript{(16)}.

%%K t01-19-4 p
\VUblancAlinea
Sur l'armoro esas : \VUgras{sfero} \textsuperscript{(14)}, \VUgras{kono} e \VUgras{ter}\cc
\VUgras{globo} \textsuperscript{(13)}.

%%K t01-19-5 p
\VUblancAlinea
Super la tabeli pendas \VUgras{mapo} reprezentanta\nl
la kin parti dil mondo : \VUgras{Europa} \textsuperscript{(g)}, \VUgras{Azia} \textsuperscript{(e)},\nl
\VUgras{Afrika} \textsuperscript{(f)}, \VUgras{Amerika} \textsuperscript{(ab)} ed \VUgras{Oceania} \textsuperscript{(h)} e la du\nl
granda Oceani, la \VUgras{Pacifiko} \textsuperscript{(c)} e l'\VUgras{Atlantiko} \textsuperscript{(d)}.

\VUsaut{6.0mm}
\VUfilet{22mm}
\end{VUpage}
